% Persona 3: Administrador del prototipo
\begingroup
\fontsize{7.5}{9.7}\selectfont
\ifdefined\pcolw\else\newlength{\pcolw}\fi
\ifdefined\pfootw\else\newlength{\pfootw}\fi
\definecolor{PInk}{RGB}{15,118,110}
\definecolor{PTeal}{RGB}{13,148,136}
\definecolor{PBlue}{RGB}{37,99,235}
\definecolor{PSlate}{RGB}{71,85,105}
\definecolor{PEmr}{RGB}{5,150,105}
\definecolor{PAmb}{RGB}{217,119,6}
\definecolor{PRed}{RGB}{220,38,38}
\definecolor{PVio}{RGB}{124,58,237}
\definecolor{PInd}{RGB}{67,56,202}
\definecolor{PSky}{RGB}{2,132,199}
\definecolor{PPlum}{RGB}{190,24,93}
\newcommand{\PDot}[1]{%
  \tikz[baseline=-0.42ex]\fill[#1] (0,0) circle (1.85pt);%
  \hspace{0.48em}%
}
\newcommand{\PTit}[1]{\makebox[\linewidth][c]{\sffamily\bfseries #1}}
\tikzset{
  pchip/.style={
    sharp corners,
    rounded corners=7.7pt,
    inner xsep=6.8pt,
    inner ysep=0pt,
    minimum height=15.4pt,
    text height=0.68em,
    text depth=0.22em,
    font=\sffamily\bfseries\scriptsize,
    draw=none,
    line width=0pt,
    outer sep=0pt,
    align=center,
    anchor=base,
  },
}
\newcommand{\PPills}[6]{%
  \begin{tikzpicture}[baseline=0.22em, rounded corners=0pt, sharp corners]
    \node[pchip, fill=#1, text=#2] (p1) {#3};
    \node[pchip, fill=#4, text=#5, base right=4.2pt of p1] (p2) {#6};
  \end{tikzpicture}%
}
\newcommand{\PCaja}[4][south]{%
  \noindent
  \begin{tikzpicture}[baseline=(current bounding box.#1)]
    \begin{scope}[shift={(0.4pt,-1.8pt)}]
      \node[
        fill=black,
        opacity=0.18,
        text opacity=0,
        rounded corners=13pt,
        inner xsep=8.4pt,
        inner ysep=5.8pt,
        outer sep=0pt,
        align=left,
        font=\fontsize{7.5}{9.7}\selectfont\sffamily,
        text width=\dimexpr#2-16.8pt\relax,
      ] {#4};
    \end{scope}
    \node[
      fill=#3!9,
      rounded corners=13pt,
      inner xsep=8.4pt,
      inner ysep=5.8pt,
      outer sep=0pt,
      text=black,
      align=left,
      font=\fontsize{7.5}{9.7}\selectfont\sffamily,
      text width=\dimexpr#2-16.8pt\relax,
    ] (pcard) {#4};
    \fill[#3, rounded corners=2.1pt]
      ([xshift=2.2pt,yshift=5.2pt]pcard.south west)
      rectangle
      ([xshift=5.7pt,yshift=-5.2pt]pcard.north west);
  \end{tikzpicture}%
}

\begin{tikzpicture}
  \node[
    rounded corners=16pt,
    inner sep=0pt,
    outer sep=0pt,
    fill=black!3.5,
  ] {%
    \begin{minipage}{0.988\linewidth}
      \begin{tikzpicture}
        \node[
          fill=PInk,
          rounded corners=12pt,
          inner sep=0pt,
          outer sep=0pt,
          text width=\linewidth-1.4pt,
        ] {%
          \begin{minipage}{\linewidth}
            \vspace{7pt}
            \hspace{9pt}%
            \begin{minipage}{0.60\linewidth}
              \begin{tikzpicture}[baseline=-0.55ex]
                \fill[white] (0,0) circle (0.34cm);
                \node[font=\sffamily\bfseries\scriptsize, color=PInk] at (0,0) {AM};
              \end{tikzpicture}%
              \hspace{6pt}%
              \raisebox{0.15ex}{\textcolor{white}{\sffamily\bfseries\Large Andr\'es Molina}}\\[3pt]
              \hspace{1pt}\textcolor{white}{\sffamily\scriptsize Mantener el acceso ordenado para que el equipo trabaje sin fricci\'on}\\[2pt]
              \textcolor{white!88}{\sffamily\tiny Administrador del prototipo $\cdot$ ART $\cdot$ persona de soporte}
            \end{minipage}%
            \hfill
            \begin{minipage}{0.34\linewidth}
              \raggedleft
              \PPills{white}{PInk}{Admin}{PSlate}{white}{Cuentas y roles}\\[5pt]
              \textcolor{white!88}{\sffamily\tiny Uso puntual, cr\'{\i}tico}
            \end{minipage}%
            \hspace{9pt}%
            \vspace{7pt}
          \end{minipage}%
        };
      \end{tikzpicture}

      \vspace{8pt}
      \raggedright
      \begin{minipage}{\dimexpr\linewidth-22pt\relax}
        \setlength{\pcolw}{\dimexpr(\linewidth-16pt)/3\relax}
        \setlength{\pfootw}{\dimexpr(\linewidth-8pt)/2\relax}

        \PCaja{\linewidth}{PTeal}{%
          {\color{PTeal}\PTit{Qui\'en es y c\'omo trabaja}}\\[4pt]
          T\'ecnico de sistemas interno de la ART, responsable de la configuraci\'on inicial y del mantenimiento operativo de Prevex.
          Conoce la estructura de roles (Preventor, Coordinador, Administrador) pero no opera en campo ni arma el plan de visitas.
          Su trabajo es que el equipo entre, que quien sale del \'area no siga entrando, y que el historial de visitas no se pierda al dar de baja una cuenta.
          Cuando hace falta, exige MFA y manda la invitaci\'on por correo para que el alta no dependa de contrase\~nas compartidas.
        }

        \vspace{9pt}
        \noindent
        \begin{minipage}[t]{\pcolw}
          \PCaja[north]{\pcolw}{PBlue}{%
            {\color{PBlue}\PTit{Contexto de uso}}\\[4pt]
            \PDot{PBlue}D\'onde: panel de usuarios y configuraci\'on del prototipo.\\[2.6pt]
            \PDot{PBlue}Cu\'ando: altas al incorporar preventores; bajas al rotar el equipo.\\[2.6pt]
            \PDot{PBlue}Por qu\'e: sin acceso ordenado el ranking y la agenda no sirven.\\[2.6pt]
            \PDot{PBlue}Competencia: roles y cuentas; no el modelo predictivo.
          }
        \end{minipage}%
        \hspace{8pt}%
        \begin{minipage}[t]{\pcolw}
          \PCaja[north]{\pcolw}{PSlate}{%
            {\color{PSlate}\PTit{Herramientas de hoy}}\\[4pt]
            \PDot{PSlate}Altas y bajas manuales o por tickets de sistemas.\\[2.6pt]
            \PDot{PSlate}Cuentas compartidas o reutilizadas entre turnos.\\[2.6pt]
            \PDot{PSlate}Poco control de qui\'en dej\'o de operar en campo.\\[2.6pt]
            \PDot{PSlate}Sin MFA ni invitaci\'on trazable por correo.
          }
        \end{minipage}%
        \hspace{8pt}%
        \begin{minipage}[t]{\pcolw}
          \PCaja[north]{\pcolw}{PEmr}{%
            {\color{PEmr}\PTit{Motivaciones}}\\[4pt]
            \PDot{PEmr}Que solo usuarios activos entren a Prevex.\\[2.6pt]
            \PDot{PEmr}Conservar el historial de visitas al dar de baja.\\[2.6pt]
            \PDot{PEmr}Evitar contrase\~nas compartidas en el equipo.\\[2.6pt]
            \PDot{PEmr}Poder exigir MFA cuando la pol\'{\i}tica lo pida.
          }
        \end{minipage}

        \vspace{9pt}
        \noindent
        \begin{minipage}[t]{\pcolw}
          \PCaja[north]{\pcolw}{PAmb}{%
            {\color{PAmb}\PTit{Objetivos}}\\[4pt]
            \PDot{PAmb}Alta, edici\'on y baja l\'ogica de usuarios (HU-10 / RF-10).\\[2.6pt]
            \PDot{PAmb}Asignar roles Preventor, Coordinador o Administrador.\\[2.6pt]
            \PDot{PAmb}Invitar por correo y exigir MFA cuando corresponda (HU-16 / RF-16).\\[2.6pt]
            \PDot{PAmb}Auditar que inactivos no puedan autenticarse.
          }
        \end{minipage}%
        \hspace{8pt}%
        \begin{minipage}[t]{\pcolw}
          \PCaja[north]{\pcolw}{PRed}{%
            {\color{PRed}\PTit{Frustraciones}}\\[4pt]
            \PDot{PRed}Procesos manuales para dar de baja usuarios.\\[2.6pt]
            \PDot{PRed}Cuentas que se reasignan y pierden trazabilidad.\\[2.6pt]
            \PDot{PRed}Borrar un usuario y arrastrar el historial de visitas.\\[2.6pt]
            \PDot{PRed}Altas lentas si hay que inventar y comunicar claves.
          }
        \end{minipage}%
        \hspace{8pt}%
        \begin{minipage}[t]{\pcolw}
          \PCaja[north]{\pcolw}{PVio}{%
            {\color{PVio}\PTit{H\'abitos}}\\[4pt]
            \PDot{PVio}Prefiere baja l\'ogica a borrar filas.\\[2.6pt]
            \PDot{PVio}Revisa roles antes de habilitar acceso al ranking.\\[2.6pt]
            \PDot{PVio}Usa invitaci\'on por correo en altas nuevas.\\[2.6pt]
            \PDot{PVio}No opera la agenda de campo: deja eso a Preventor y Coordinaci\'on.
          }
        \end{minipage}

        \vspace{10pt}
        \PCaja{\linewidth}{PInd}{%
          {\color{PInd}\PTit{Escenario}}\\[4pt]
          Entra un preventor nuevo al \'area: Andr\'es crea la cuenta, asigna rol Preventor y env\'{\i}a la invitaci\'on por correo.
          Si la pol\'{\i}tica lo exige, marca MFA obligatorio antes del primer acceso completo.
          Meses despu\'es esa persona rota de sector: la da de baja l\'ogica.
          Las visitas que ya carg\'o siguen visibles en el historial del establecimiento; solo deja de autenticarse.
        }

        \vspace{9pt}
        \noindent
        \begin{minipage}[t]{\pfootw}
          \PCaja[north]{\pfootw}{PSky}{%
            {\color{PSky}\PTit{Qu\'e necesita de Prevex}}\\[4pt]
            \PDot{PSky}ABM de usuarios con baja l\'ogica (HU-10 / RF-10).\\[2.6pt]
            \PDot{PSky}Invitaci\'on y recuperaci\'on por correo (HU-16 / RF-16).\\[2.6pt]
            \PDot{PSky}MFA TOTP opcional u obligatorio por usuario (RNF-01).\\[2.6pt]
            \PDot{PSky}Que el historial de visitas no dependa de que la cuenta siga activa.
          }
        \end{minipage}%
        \hspace{8pt}%
        \begin{minipage}[t]{\pfootw}
          \PCaja[north]{\pfootw}{PPlum}{%
            {\color{PEmr}\PTit{\'Exito}}\\[3pt]
            Alta, edici\'on o baja l\'ogica en menos de dos minutos, sin perder visitas asociadas ni dejar cuentas hu\'erfanas.\\[5pt]
            {\color{PPlum}\PTit{Temor}}\\[3pt]
            Que un ex integrante siga entrando; que borrar una cuenta borre el rastro operativo del establecimiento.
          }
        \end{minipage}

        \vspace{4pt}
      \end{minipage}
    \end{minipage}%
  };
\end{tikzpicture}
\endgroup
